% =====================================================================
% SEC_03 — Fundamentação Teórica
% =====================================================================

\section*{3. Fundamentação Teórica}
\addcontentsline{toc}{section}{3. Fundamentação Teórica}

\nivel{0} Esta seção constrói a base conceitual que sustenta tanto o estado atual
quanto a proposta. Percorremos, em ordem crescente de complexidade: (1) como
ranqueamos documentos (BM25); (2) como acoplamos recuperação e geração (RAG);
(3) como melhoramos o RAG com raciocínio e correção; (4) como diversificamos
(MMR); e (5) que ferramenta matemática usamos para tornar a diversificação
determinística (Recamán). Cada tópico começa no nível 0 (intuição) e evolui até o
rigor formal.

\subsection*{3.1 Ranqueamento clássico: BM25}

\nivel{0} Quando você digita uma busca, o sistema precisa decidir a ordem dos
resultados. A intuição do BM25 é: um documento é mais relevante se contém os
termos da busca com alta frequência \emph{relativa}, e termos que aparecem em
poucos documentos são mais "informativos" do que termos que aparecem em todos.

\begin{intuicao}
Se você busca "baleia" e um documento usa a palavra "baleia" muitas vezes, ele
provavelmente é relevante. Mas a palavra "o" aparece em todo documento e, portanto,
diz pouco. O BM25 equilibra "quantas vezes o termo aparece" com "quão raro o termo
é no corpus".
\end{intuicao}

\nivel{2} O \emph{Best Matching 25} (BM25) é um dos modelos de ranqueamento mais
estabelecidos, formalizado no quadro da teoria da relevância probabilística
\cite{RobertsonZaragoza2009BM25}. Para um documento $D$ e uma consulta $Q$
contendo os termos $q_1,\dots,q_n$, o escore é dado por:

\begin{equation}
\label{eq:bm25}
\text{BM25}(D,Q)\;=\;\sum_{i=1}^{n}
\text{IDF}(q_i)\;\cdot\;
\frac{f(q_i,D)\,(k_1+1)}{f(q_i,D)+k_1\,\bigl(1-b+b\,\tfrac{|D|}{\text{avgdl}}\bigr)}
\end{equation}

onde $f(q_i,D)$ é a frequência do termo $q_i$ em $D$, $|D|$ é o comprimento do
documento, $\text{avgdl}$ é o comprimento médio dos documentos do corpus, e
$k_1,b$ são parâmetros de saturação e de normalização por comprimento. O termo de
frequência inversa de documento $\text{IDF}(q_i)$ é função do número de documentos
que contêm $q_i$. Este modelo estabelece o \emph{baseline} lexical do
ranqueamento no ecossistema.

\nivel{1} Para este manual, o que importa reter é: o BM25 produz uma \emph{ordem}
dos documentos. É essa ordem que a proposta de diversificação irá
\emph{re-balancear} em um estágio posterior, sem refazer o ranqueamento lexical.

\subsection*{3.2 RAG canônico}

\nivel{0} RAG significa "gerar texto a partir de documentos recuperados". Em vez
de um modelo de linguagem responder "de memória", ele primeiro \emph{recupera}
fontes relevantes e depois escreve a resposta \emph{apoiado} nessas fontes. Isso
reduz alucinações e permite atualizar o conhecimento sem retreinar o modelo.

\begin{intuicao}
Duas estratégias: (a) um aluno responde de cabeça --- rápido, mas pode errar; (b)
um aluno abre os livros antes de responder --- mais confiável. O RAG é a
estratégia (b): o modelo consulta a "biblioteca" (recuperador) e então responde
(gerador).
\end{intuicao}

\nivel{2} O paradigma RAG combina um recuperador (retriever) com um gerador
sequência-a-sequência. O arcabouço original \cite{Lewis2020RAG} condiciona a
geração nos documentos recuperados $z \in \mathcal{Z}$, modelando
$p(y \mid x) \approx \sum_{z \in \mathcal{Z}} p_\eta(z \mid x)\; p_\theta(y \mid x, z)$,
onde $\eta$ parametriza o recuperador e $\theta$ o gerador. A recuperação densa
por interação tardia (\emph{late interaction}), como em ColBERT
\cite{KhattabZaharia2020ColBERT}, permite forte poder expressivo com custo de
inferência reduzido.

\nivel{1} Conceitualmente, o RAG tem dois fatores que importam para este manual:
o \emph{recuperador} (decide \emph{o que} entra no contexto) e o \emph{gerador}
(decide \emph{como} responder). A diversificação atua no recuperador, mudando o
\emph{quê} entra --- não o \emph{como} é respondido.

\subsection*{3.3 Recuperação aumentada e raciocínio}

\nivel{0} O RAG simples recupera uma vez e responde. Mas perguntas complexas
precisam de \emph{vários passos}: recuperar, raciocinar, recuperar de novo. Além
disso, o RAG pode "quebrar" se os documentos recuperados forem ruins --- existem
mecanismos para detectar e corrigir isso.

\nivel{2} Trabalhos subsequentes demonstraram que a qualidade do RAG depende
criticamente da \emph{posição} dos documentos relevantes no contexto longo
\cite{Liu2024LostMiddle}, e que estratégias iterativas de recuperação combinada
com raciocínio (\emph{retrieve-then-read}, \emph{interleave}) melhoram a precisão
em perguntas de múltiplos passos \cite{IzacardGrave2021LeveragingPassages,
Trivedi2023Interleaving}. O RAG auto-adaptativo roteia a intensidade da
recuperação conforme a complexidade da consulta \cite{Jeong2024AdaptiveRAG}.
Abordagens baseadas em grafos, como o \emph{GraphRAG}, organizam o corpus em
comunidades temáticas para sumarização focada na consulta \cite{Edge2024GraphRAG}.
Métodos de correção pós-recuperação tratam a degradação quando a recuperação
falha \cite{Yan2024CRAG}, e há evidências de que o uso da recuperação em tempo de
inferência (\emph{in-context}) requer ajuste fino eficiente
\cite{Ram2023RetrievalAugmented}. Alternativas \emph{prompt-only}, como o
raciocínio com recuperação sobre passos de raciocínio, buscam fidelidade sem novo
treinamento \cite{He2023RethinkingRetrieval}.

\begin{pontochave}
O RAG é um campo em evolução. A proposta deste manual (diversificação por Recamán)
é \emph{ortogonal} a esses avanços: ela não compete com raciocínio iterativo,
correção ou grafos --- atua como um estágio de pós-processamento do conjunto
recuperado, podendo ser combinada com todos eles.
\end{pontochave}

\subsection*{3.4 Diversificação: MMR}

\nivel{0} Se o RAG recupera documentos, a \emph{diversificação} decide como evitar
redundância: entre dois documentos igualmente relevantes, prefira aquele que
\emph{adiciona informação nova} em vez de repetir o que já foi dito.

\begin{intuicao}
Ao montar uma comissão para debater um tema, você não convida dez integrantes
iguais. Você equilibra \emph{relevância} (pessoas entendidas) com
\emph{diversidade} (pessoas com pontos de vista distintos). O MMR é esse
"recrutador".
\end{intuicao}

\nivel{2} A \emph{Maximal Marginal Relevance} (MMR) equilibra relevância e
novidade \cite{CarbonellGoldstein1998MMR}. Dado um conjunto de documentos $D$ já
selecionado e o restante $R$, o próximo documento $d$ é escolhido por:

\begin{equation}
\label{eq:mmr}
\text{MMR} \;=\; \arg\max_{d \in R}\Bigl[
\lambda\,\text{Sim}_1(d,Q)\;-\;(1-\lambda)\,
\max_{d_j \in D}\text{Sim}_2(d,d_j)
\Bigr]
\end{equation}

O parâmetro $\lambda \in [0,1]$ controla o balanço entre relevância à consulta
($\text{Sim}_1$) e redundância em relação aos já selecionados ($\text{Sim}_2$).
MMR é a inspiração central da proposta de diversificação deste manual.

\nivel{2} Um ponto crítico para a nossa proposta: o MMR depende de uma medida de
similaridade $\text{Sim}_2$ e de um parâmetro $\lambda$. Essa dependência pode
introduzir \emph{não-determinismo} (variação conforme a implementação da
similaridade) e sensibilidade a $\lambda$. A proposta deste manual busca
\emph{uma alternativa determinística} -- ainda inspirada no espírito de MMR, mas
baseada em uma sequência canônica (Recamán) com cobertura determinística do
ranking.

\subsection*{3.5 A sequência de Recamán}

\nivel{0} A sequência de Recamán é uma "conta" matemática simples com um "pulos"
curiosos. Começa em 0. Em cada passo, você tenta \emph{subtrair} o número do passo;
se o resultado der negativo ou já tiver aparecido, você \emph{soma}. O resultado é
uma sequência que salta entre números pequenos e grandes, criando uma "cobertura"
irregular mas **totalmente previsível** -- determinística.

\begin{intuicao}
Imagine escolher cadeiras em uma sala muito grande, sempre tentando andar para
trás; se não der, andando para frente. A sequência de Recamán produz um padrão de
saltos que "espalha" os números de forma que nenhum valor se repete precocemente.
Para a diversificação, esse "espalhamento" é exatamente o que queremos: evitar
ficar preso a um único canto do ranking.
\end{intuicao}

\nivel{2} A sequência de Recamán é definida por $a_0 = 0$ e, para $n \geq 1$:

\begin{equation}
\label{eq:recaman}
a_n = \begin{cases}
a_{n-1} - n, & \text{se } a_{n-1} - n > 0 \text{ e } a_{n-1} - n \notin \{a_0,\dots,a_{n-1}\},\\
a_{n-1} + n, & \text{caso contrário}.
\end{cases}
\end{equation}

A sequência produz uma permutação dos inteiros não negativos e exibe propriedades
misteriosas ainda em estudo; "primos da sequência" foram descritos em
\cite{Alekseyev2022RecamanCousins}. A proposta deste manual usa a sequência como
fonte determinística de \emph{offsets} de diversificação, conferindo
reprodutibilidade total ao estágio de diversificação.

\begin{pontochave}
A ponte conceitual está pronta: o MMR inspira \emph{por que} diversificar
(Eq.~\ref{eq:mmr}); a sequência de Recamán oferece \emph{como} diversificar de
forma determinística (Eq.~\ref{eq:recaman}). A próxima seção mostra o que já está
implementado no ecossistema.
\end{pontochave}
